% ---------------------------------------------------------------------------
% NeurIPS / ICML Conference Paper Template
%
% Usage:
%   1. Download official neurips_2026.sty from: https://neurips.cc/Conferences/2026/PaperInformation/StyleFiles
%   2. Place neurips_2026.sty in this directory
%   3. Compile: pdflatex template.tex && bibtex template && pdflatex template.tex x2
%
% Conference: NeurIPS (Neural Information Processing Systems) / ICML (International Conference on Machine Learning)
% Page limit: 9 pages (NeurIPS) / 8 pages (ICML) + references + appendices
% ---------------------------------------------------------------------------

\documentclass{article}

% --- UNCOMMENT when neurips_2026.sty is available ---
% \usepackage[final]{neurips_2026}
% \usepackage[utf8]{inputenc}
% \usepackage[T1]{fontenc}
% \usepackage{hyperref}
% \usepackage{url}
% \usepackage{booktabs}
% \usepackage{amsfonts,amsmath,amssymb}
% \usepackage{nicefrac}
% \usepackage{microtype}
% \usepackage{xcolor}

\usepackage[utf8]{inputenc}
\usepackage[T1]{fontenc}
\usepackage{graphicx}
\usepackage{amsmath,amssymb}
\usepackage{booktabs}
\usepackage{multirow}
\usepackage[pagebackref,breaklinks,colorlinks,citecolor=blue]{hyperref}
\usepackage[margin=1in]{geometry}  % Remove when using neurips_2026.sty

% --- Paper ID (anonymous submission) ---
% \neuripsfinalcopy

\title{Your Paper Title Here}

\author{
  Anonymous Author(s) \\
  Anonymous Institution \\
  \texttt{anonymous@institution.edu}
}

\begin{document}

\maketitle

% ===========================================================================
% ABSTRACT
% ===========================================================================
\begin{abstract}
  Your abstract here. NeurIPS typically allows slightly longer abstracts than CVPR.
  Follow 2-6-2 structure: 2 background/gap, 6 method (First/Second/Finally),
  2 results/implication. Target: $\sim$210 words.
\end{abstract}

% ===========================================================================
% 1. INTRODUCTION
% ===========================================================================
\section{Introduction}

% Standard 5-paragraph structure with NeurIPS flavor:
% - More emphasis on algorithmic/optimization novelties
% - Theory/analysis contributions highlighted more than CVPR
% - Broader impact statement often required

Your introduction text here.

% ===========================================================================
% 2. RELATED WORK
% ===========================================================================
\section{Related Work}

% ===========================================================================
% 3. METHOD
% ===========================================================================
\section{Method}

% Opening paragraph (3-5 sentences): overall framework → what each innovation
% solves → how they collaborate. No "Overview" or "Preliminaries" subsection.
% Reference the architecture diagram.

\begin{figure}[t]
  \centering
  \fbox{\parbox{0.8\linewidth}{\centering
  \vspace{2cm}
  {\large [Architecture Diagram]}
  \vspace{2cm}}}
  \caption{Overview of the proposed method.}
  \label{fig:architecture}
\end{figure}

\subsection{<Innovation 1 Name>}
\label{sec:method-innov1}

% Motivation → Design → Formula/Mechanism → Effect

\subsection{<Innovation 2 Name>}
\label{sec:method-innov2}

% Motivation → Design → Formula/Mechanism → Effect

\subsection{<Innovation 3 Name>}
\label{sec:method-innov3}

% Motivation → Design → Formula/Mechanism → Effect

\subsection{Implementation Details}
\label{sec:method-impl}

% Optimizer, LR schedule, batch size, data augmentation, etc.

% ===========================================================================
% 4. THEORETICAL ANALYSIS (if applicable)
% ===========================================================================
% \section{Theoretical Analysis}
% \label{sec:theory}

% ===========================================================================
% 5. EXPERIMENTS
% ===========================================================================
\section{Experiments}
\label{sec:experiments}

\subsection{Datasets and Metrics}
\label{sec:exp-data}

\subsection{Experimental Setup}
\label{sec:exp-setup}

\subsection{Comparison with State-of-the-Art and Classical Methods}
\label{sec:exp-main}

\subsection{Ablation Study}
\label{sec:exp-ablation}

% Numbered list 1), 2), 3)... analyzing each component / design choice

\subsection{Visualization}
\label{sec:exp-visual}

% Feature maps, attention maps, t-SNE, error cases, etc.

% ===========================================================================
% 6. CONCLUSION
% ===========================================================================
\section{Conclusion}
\label{sec:conclusion}

% Paragraph 1: Full summary (problem → method → key results with numbers)
% Paragraph 2: Limitations + future work (specific limitations + actionable directions)

% ===========================================================================
% BROADER IMPACT (required for NeurIPS)
% ===========================================================================
\section*{Broader Impact}

Discuss potential societal consequences of this work, both positive and negative.

% ===========================================================================
% ACKNOWLEDGMENTS
% ===========================================================================
\section*{Acknowledgments}

% ===========================================================================
% REFERENCES
% ===========================================================================
{
  \small
  \bibliographystyle{plain}
  \bibliography{references}
}

% ===========================================================================
% APPENDIX (optional)
% ===========================================================================
\appendix
\section{Appendix: Additional Results}
\label{sec:appendix}

% Additional experiments, proofs, implementation details.

\end{document}
